\section{Discussion}
\label{sec:discussion}

\subsection{Points forts}

\begin{itemize}
  \item Le pipeline principal (collecte, indexation, RAG CLI) est \textbf{fonctionnel de bout en
        bout} et produit des r\'eponses ancr\'ees dans les sources Wikipedia du domaine NLP/ML.
  \item Les r\'esultats d'\'evaluation montrent un \textbf{recall@3 moyen de 0.985} et une
        \textbf{precision@3 moyenne de 0.922} : le retriever bas\'e sur \texttt{all-MiniLM-L6-v2}
        retrouve quasi syst\'ematiquement le bon document, m\^eme pour des formulations diff\'erentes
        du texte source.
  \item La \textbf{fid\'elit\'e moyenne de 0.887} et le prompt restrictif limitent les
        hallucinations dans la majorit\'e des cas observ\'es.
  \item Le \textbf{temps de r\'eponse moyen de 0.81~s} rend le syst\`eme utilisable de mani\`ere
        interactive via le CLI (et, une fois valid\'ee, via Streamlit).
  \item Le d\'ep\^ot GitHub documente chaque phase (\texttt{PHASES.md}) et centralise les
        param\`etres dans \texttt{config.py}, ce qui facilite la reproductibilit\'e.
\end{itemize}

\subsection{Limites et travail en cours}

\begin{itemize}
  \item \textbf{Documents voisins} : pour des sujets proches (FastText / Word2vec, sentence
        embedding / BERT), le retriever renvoie parfois des chunks pertinents s\'emantiquement mais
        issus d'un document diff\'erent de celui marqu\'e comme ``attendu'', ce qui p\'enalise la
        precision@3 sans invalider la qualit\'e de la r\'eponse.
  \item \textbf{M\'etrique de fid\'elit\'e lexicale} : elle ne d\'etecte ni les reformulations
        correctes (synonymes) ni les hallucinations subtiles r\'eutilisant le vocabulaire du
        contexte. Une \'evaluation par LLM-as-judge serait plus fine mais plus co\^uteuse.
  \item \textbf{Corpus en anglais} : le syst\`eme r\'epond en anglais et ne couvre que le contenu
        des articles t\'el\'echarg\'es ; il ne refl\`ete pas les avanc\'ees post\'erieures au
        t\'el\'echargement.
  \item \textbf{Interface Streamlit} : impl\'ement\'ee avec des fonctionnalit\'es avanc\'es
        (\'evaluation int\'egr\'ee, gestion du corpus), mais sa validation utilisateur finale et
        la pr\'eparation de la d\'emonstration orale restent en cours au moment de la r\'edaction.
  \item \textbf{Pas de re-ranking} : la recherche vectorielle seule ne desambiguise pas les
        articles tr\`es proches ; un cross-encoder ou un filtrage par m\'etadonn\'ees pourrait
        am\'eliorer la precision.
\end{itemize}

\subsection{Alignement avec le sujet du projet}

Le projet respecte les exigences principales du sujet : corpus $\geq$30 documents, index
vectoriel, pipeline RAG avec LLM, interface (CLI + Streamlit), jeu de test $\geq$20 questions,
m\'etriques demand\'ees et graphiques de performance. Les parties encore en cours (tests Streamlit,
README final, d\'emo) sont explicitement signal\'ees et n'invalident pas les r\'esultats
d'\'evaluation obtenus sur le pipeline principal.
